\section{Use a reconstructed picture as the next reference}
The completed macroblock emitter knows how to serialize a selected vector and
six blocks. It does not know where the blocks belong in a picture. The picture
controller supplies that geometry and the order of operations. Its input is one
complete raw QCIF picture, plus the preceding locally reconstructed picture for
a P-picture. Its two outputs are an encoded picture and a new reconstruction.
Neither output is the original raw input with a different name.

Use raster order: nine macroblock rows, with eleven macroblocks in each row.
For row $r$ and column $c$, the luma origin is $(16c,16r)$ and the chroma origin
is $(8c,8r)$. Four luma blocks cover the sixteen-by-sixteen region; one block
covers each chroma plane. Together the $99\times6$ blocks cover exactly
$99\times6\times64=38\,016$ active samples. This equality also checks that a
controller has neither omitted nor duplicated any output samples.

An I-picture selects INTRA for every macroblock in this controller. For a
P-picture, perform the earlier bounded search against the \emph{reconstructed}
reference, form six prediction records, and encode all six INTER blocks.
This is a stated teaching policy, not a requirement that every H.263 encoder
must use INTER for every P-picture macroblock. An independent encoder can select
INTRA inside a P-picture. Our earlier macroblock entry supports that syntax;
this controller deliberately develops one mode-selection policy at a time.

For each completed block, retain its selected levels for macroblock syntax and
copy its reconstructed pixels into the private destination picture. Copy each
row to the plane's stride, rather than placing every sixty-four-byte block
consecutively in the image. Consecutive block records are a working layout;
consecutive picture samples are a different layout.

\section{Make group boundaries part of the predictor design}
The general baseline vector predictor is the component-wise median of the left,
above and above-right candidates, with the border rules of clause 6.1.1 applied
in order. INTRA and skipped neighbours contribute zero. An unavailable left
candidate contributes zero. At the top of the picture, or the top of a GOB whose
header is nonempty, the two upper candidates become the left candidate. An
above-right position beyond the picture's right edge then becomes zero.

A QCIF group of blocks contains one macroblock row. This controller emits a
nonempty GOB header for rows one through eight. Row zero has no separate GOB
header: its picture header already establishes the first group. Each row therefore
starts with a zero predictor. Later macroblocks in that row use the actual vector
of the preceding macroblock. At the right edge the median of left, left and zero
is still left. This row-based state follows the general rule; it is not a
replacement predictor that a decoder must be told to understand.

Before a GOB start code, insert fewer than eight zero stuffing bits to reach a
byte boundary. Emit these fields in order:
\begin{center}
\tagpdfsetup{table/header-rows={1}}
\begin{tabular}{lrl}
\textbf{Field} & \textbf{Bits} & \textbf{Value here}\\
GBSC & 17 & One, as a seventeen-bit code\\
GN & 5 & Current row, one through eight\\
GFID & 2 & The supplied picture identity\\
GQUANT & 5 & The picture quantizer, unchanged
\end{tabular}
\end{center}
GFID is the same in every group of one picture. Without extended PTYPE, it stays
the same when PTYPE stays the same and changes when PTYPE changes. A clip with
one I-picture followed by P-pictures can use zero for the first picture and one
for every later picture. It must not alternate GFID merely because a new picture
has begun. The caller tracks that history and supplies the identity.

After emitting a macroblock, retain its selected actual vector for the next
column. The syntax request uses twice the full-pixel displacement. A skip is
possible only with actual vector zero and no residual, so retaining that vector
also records the correct skipped-neighbour candidate. Reset both components when
starting the next explicit GOB row.

\section{Stage a complete picture before publication}
The destination reconstruction must not become the reference while its lower
rows are still uninitialized. Use separate old-reference and new-picture storage.
Likewise, a syntax failure near the last macroblock must not leave a short picture
advertised as a completed output.

The controller receives a caller-owned workspace with two private regions:
\begin{center}
\tagpdfsetup{table/header-rows={1}}
\begin{tabular}{rrl}
\textbf{Offset} & \textbf{Bytes} & \textbf{Purpose}\\
0 & 131\,072 & Encoded-picture staging arena\\
131\,072 & 38\,016 & New reconstructed Y, Cb and Cr planes
\end{tabular}
\end{center}
Its total is 169\,088 bytes. This workspace may change during a failed call;
it is not either published result. The old reference remains read-only throughout.
Stack storage holds the search, prediction, block and syntax records, not another
whole picture. Supply at least 65\,536 bytes of private stack for this composition.

The largest supported INTER macroblock has 8487 payload bits. Ninety-nine such
payloads, fifty picture-header bits, eight group headers with at most seven
stuffing bits each, and final byte alignment fit within 105\,070 bytes. The
131\,072-byte staging arena therefore admits the complete supported picture.
This bound concerns one picture; it is not a limit on the whole video file.

The stage starts at bit count zero. There is no need to clear all 131\,072 bytes
first: every published bit is supplied by a field, a macroblock or explicit zero
alignment. Earlier bit writers preserve bytes outside their touched windows, but
those bytes are never published as part of the picture. This avoids a repeated
large initialization while keeping output deterministic.

After all dependencies succeed, append zero alignment to the next byte boundary.
Compare the complete staged byte count with caller capacity. Only then copy the
encoded bytes and all 38\,016 reconstructed samples into the caller's two output
areas. Commit the caller's used-bit count last. These final mapped-memory copies
have no fallible children. An output-file write is a later operation and needs its
own acceptance check before the outer clip loop swaps its reference buffers.
